% Source: https://artofproblemsolving.com/community/c3h1991373_directly_and_inversely_proportional
% Stable identity: aops-topic-1991373-q01
From the first chain, cross-multiplying,
\[
x(y+z)=y(x+y)\ \Longrightarrow\ xz=y^2,\qquad y(z+x)=z(y+z)\ \Longrightarrow\ xy=z^2 .
\]
Dividing the two results, $\dfrac {z}{y}=\dfrac{y^2}{z^2}$, so $y^3=z^3$ and $y=z$; then $xz=z^2$ gives $x=z$. So $x=y=z$.

From the second chain, $y(y+z)=z(z+t)$ with $y=z$ gives $2z^2=z^2+zt$, so $t=z$ (and then $t(t+y)=2z^2$ as well). Hence all four angles are equal and, since they fill the full angle around $O$,
\[
x=y=z=t=\boxed{90^\circ}.
\]

\begin{asy}
usepackage("amsmath");
size(3.5cm);
defaultpen(fontsize(10pt));
pair O=(0,0);
draw((-1.2,0)--(1.2,0),linewidth(0.9)); draw((0,-1.2)--(0,1.2),linewidth(0.9));
dot(O); label("$O$",O,SW);
label("$x$",(0.45,0.45)); label("$y$",(-0.45,0.45)); label("$z$",(-0.45,-0.45)); label("$t$",(0.45,-0.45));
draw((0.13,0)--(0.13,0.13)--(0,0.13));
label("\mbox{}",(min(currentpicture,user=true)+max(currentpicture,user=true))/2,black); // keeps the text after the figure black when LaTeX typesets the labels inline
\end{asy}

So the four angles are formed by two perpendicular lines through $O$.
